% ========================================
% 术语表
% ========================================

\chapter*{术语表}
\addcontentsline{toc}{chapter}{术语表}

\begin{description}[leftmargin=4cm, style=nextline]

\item[\highlight{智慧水利}] 运用物联网、大数据、人工智能、云计算等新一代信息技术，与水利工程、水利管理、水利服务深度融合，构建具有预报、预警、预演、预案功能的智慧水利体系。

\item[\highlight{数字孪生}] 运用数字化手段，集成水利工程全要素信息，构建物理水利工程的虚拟映射，实现与物理实体的实时交互和动态优化。

\item[\highlight{微服务架构}] 一种软件架构模式，将大型应用程序拆分为多个小型、独立的服务，每个服务都有自己的业务逻辑和数据库，通过轻量级通信机制进行交互。

\item[\highlight{容器化}] 使用容器技术（如Docker）将应用程序及其依赖项打包在一起，实现应用程序的快速部署、扩展和管理。

\item[\highlight{DevOps}] 开发（Development）和运维（Operations）的组合，是一套实践和工具，旨在自动化和集成软件开发和IT运维之间的过程。

\item[\highlight{API网关}] 作为系统入口点的服务器，负责请求路由、组合和协议转换，是微服务架构中的重要组件。

\item[\highlight{负载均衡}] 将工作负载分布到多个计算资源上的技术，以优化资源使用、最大化吞吐量、减少响应时间并避免过载。

\item[\highlight{缓存}] 存储数据的硬件或软件组件，以便更快地处理将来对该数据的请求，常用技术包括Redis、Memcached等。

\item[\highlight{消息队列}] 用于在分布式系统中传递消息的中间件，实现异步通信和解耦，常用技术包括RabbitMQ、Apache Kafka等。

\item[\highlight{数据库分片}] 将大型数据库分解为更小、更快、更易于管理的部分（称为数据分片）的数据库架构模式。

\item[\highlight{高可用}] 系统在长时间内持续运行的能力，通常通过冗余、故障转移和负载分布等技术实现。

\item[\highlight{可扩展性}] 系统处理增长工作量的能力，包括水平扩展（增加更多机器）和垂直扩展（增加单机性能）。

\item[\highlight{监控告警}] 对系统运行状态进行实时监控，并在出现异常时及时发出警报的机制。

\item[\highlight{日志聚合}] 收集、存储和分析来自多个系统组件的日志数据的过程，用于故障排查和性能分析。

\item[\highlight{持续集成}] 开发实践，要求开发人员定期将代码集成到共享存储库中，每次集成都通过自动化构建来验证。

\item[\highlight{持续部署}] 软件工程方法，所有代码更改都会自动部署到生产环境，无需人工干预。

\item[\highlight{蓝绿部署}] 减少停机时间和风险的部署策略，通过维护两个相同的生产环境（蓝和绿）来实现零停机部署。

\item[\highlight{灰度发布}] 能够平滑过渡的一种发布方式，通过切换路由权重，逐步将流量从旧版本切换到新版本。

\item[\highlight{服务网格}] 专用的基础设施层，用于服务间通信，使其快速、可靠和安全，代表技术包括Istio、Linkerd等。

\item[\highlight{云原生}] 构建和运行可弹性扩展的应用的方法论，充分利用云计算的优势，包括容器化、微服务、DevOps等技术。

\item[\highlight{物联网}] 通过信息传感设备，按约定的协议，将物品与互联网相连接，进行信息交换和通信的网络。

\item[\highlight{边缘计算}] 在网络边缘进行数据处理的计算模式，减少数据传输延迟，提高响应速度。

\item[\highlight{人工智能}] 模拟人类智能的技术，包括机器学习、深度学习、自然语言处理等子领域。

\item[\highlight{大数据}] 具有4V特征（Volume体量大、Velocity速度快、Variety类型多、Value价值密度低）的数据集合。

\item[\highlight{数据湖}] 存储大量原始数据的中央存储库，数据可以是结构化、半结构化或非结构化的。

\item[\highlight{数据仓库}] 为支持管理决策而专门设计的数据存储和分析系统，通常包含历史数据。

\item[\highlight{ETL}] 数据处理过程，包括提取（Extract）、转换（Transform）和加载（Load）三个步骤。

\item[\highlight{实时计算}] 对输入数据能够在毫秒或微秒级别的时间内给出处理结果的计算模式。

\item[\highlight{流式处理}] 对连续数据流进行实时处理的计算模式，常用技术包括Apache Storm、Apache Flink等。

\item[\highlight{批处理}] 一次性处理大量数据的计算模式，常用技术包括Apache Hadoop、Apache Spark等。

\end{description}
